% GENERATED FILE. Edit canonical lessons, then run npm run generate:books.
% canonical-source-hash: fnv1a64:df6fb15ff943b7f9
% canonical-lessons: AR-W10-ayn, AR-W11-ha-and-ta-marbuta, AR-W12-maa-salama, AR-C04-maa-with, AR-C04-al-salama, AR-C04-maa-salama, AR-C04-ila-liqaa, AR-C04-practice

\chapter{Farewells}
\label{ch:ar-farewells}
\hlchaptermodality{4}

\begin{chapteropening}
\textbf{By the end of this chapter:} \emph{I can close a conversation in Arabic, choosing between maʿa s-salāma and ilā l-liqāʾ.}

The four-chapter conversation end to end -- greet, name, how-are-you, farewell -- with every -ka swapped to -ki on the second run.
\end{chapteropening}

\section[ʿayn]{\ar{ع} (ʿayn) — the eye that became an O}

\label{lesson:AR-W10-ayn}



\begin{quote}
This is the letter people think of when they think \textquotedblleft{}Arabic.\textquotedblright{} \textbf{\ar{ع}} (\emph{ʿayn}) writes a sound made \textbf{deep in the throat} — one English simply does not have. And its history holds the best surprise in the whole alphabet: this consonant is the ancestor of our vowel \textbf{O}.
\end{quote}

\subsection*{Script — Draw it}
Break it apart — \emph{an open curve, then a bowl}:

1. \textbf{the open-mouth curve} on top — like a small \textbf{c} opening to the left. 2. \textbf{the bowl below}, when \emph{ʿayn} stands alone or ends a word — a scoop underneath.

Joined in the middle of a word it flattens into a little wedge; standing alone it shows the full curve-over-bowl. It is a \textbf{connector}, so it changes coat by position like \emph{sīn} and \emph{lām}.

\subsection*{Script — A third dots-family}
You know the \textbf{bowl} family (\ar{ب}/\ar{ن}/\ar{ت}/\ar{ث}) from Lesson 4 and the \textbf{hook-and-tail} family (\ar{ح}/\ar{خ}/\ar{ج}) from Lesson 7. \emph{ʿayn} opens a \textbf{third}: its skeleton is shared with \textbf{\ar{غ}} (\emph{ghayn}), which is the \textbf{same shape plus one dot above}. Same trick, third time — by now you should expect it.

\subsection*{Script — The surprise: ʿayn $\to$ O}
\emph{ʿayn} comes from Phoenician \textbf{ʿayin}, meaning \textquotedblleft{}\textbf{eye}\textquotedblright{} — and the oldest forms were drawn as one. When the \textbf{Greeks} borrowed the Phoenician alphabet, they had \textbf{no throat-sound} like \emph{ʿ}. So they did what borrowers do with a sign they can't use: they \textbf{repurposed} it — and turned the eye-shaped letter into the vowel \textbf{omicron}, our \textbf{O}.

So the round \textbf{O} you write is a Phoenician \textbf{eye}, and Arabic's \emph{ʿayn} is the same sign that kept its original consonant job. (Compare the thread so far: \emph{ʾālep}$\to$\textbf{A}, \emph{bēt}$\to$\textbf{B}, \emph{mem}$\to$\textbf{M}, \emph{kaph}$\to$\textbf{K}, \emph{rēsh}$\to$\textbf{R} — and now \emph{ʿayin}$\to$\textbf{O}.)

\subsection*{Script — Where you've heard it}
You have been \emph{saying} this sound since Chapter 1: \textbf{\ar{السلام} \ar{عليكم}} (\emph{as-salāmu \textbf{ʿ}alaykum}) — the \emph{ʿ} in \emph{ʿalaykum} is this letter. And the farewell you'll write two lessons from now, \textbf{\ar{مع} \ar{السلامة}}, opens with it.

\subsection*{Your turn}
\begin{itemize}
\raggedright
  \item \emph{Write it:} \textbf{\ar{ع}} — the open curve, then the bowl beneath
  \item \emph{Say it:} the sound, from the throat — as in \emph{ʿalaykum}
  \item \emph{Say it:} the family — \textquotedblleft{}\emph{ʿayn} \ar{ع}, and \emph{ghayn} \ar{غ} = the same shape \textbf{plus one dot}\textquotedblright{}
  \item \emph{Say it:} \textquotedblleft{}\emph{ʿayin} = \textbf{eye} $\to$ the Greeks made it \textbf{O}\textquotedblright{}
\end{itemize}

\begin{tcolorbox}[breakable,colback=teal!4,colframe=teal!35!black,title={Before you move on}]
Describe \textbf{\ar{ع}}'s two pieces. (An \textbf{open-mouth curve} on top, a \textbf{bowl} below when isolated or final.) What did Phoenician \emph{ʿayin} mean, and what letter did the Greeks turn it into? (\textquotedblleft{}\textbf{Eye}\textquotedblright{}; they had no such sound, so they made it the vowel \textbf{O}.) Which letter shares its skeleton? (\textbf{\ar{غ}} \emph{ghayn} — one dot above.) Next: \textbf{\ar{ه}} and the feminine ending \textbf{\ar{ة}}.
\end{tcolorbox}
\section[hāʾ, tāʾ marbūṭa]{\ar{ه} and \ar{ة} — the loop, and the ending that means \textquotedblleft{}feminine\textquotedblright{}}

\label{lesson:AR-W11-ha-and-ta-marbuta}



\begin{quote}
Two related shapes. \textbf{\ar{ه}} (\emph{hāʾ}) is a plain \textquotedblleft{}h\textquotedblright{} — and the most \textbf{shape-shifting} letter you've met, wearing four genuinely different coats. Then \textbf{\ar{ة}}, a letter built \emph{out of} it, whose whole job is to say one grammatical thing: \textbf{this word is feminine}.
\end{quote}

\subsection*{Script — \ar{ه} (hāʾ) — the loop that never looks the same}
Break it apart — there is really just \emph{one piece}:

1. \textbf{the loop} — a closed round body.

But that loop \textbf{changes shape a lot by position} more than any letter so far: standing alone it's a neat circle-ish loop; at the start it opens out; in the middle it becomes a small figure-of-eight pinched in the centre; at the end it trails. Same letter, four faces — the four-coats rule from Lesson 2 taken to its extreme.

From Phoenician \textbf{hē} — the ancestor of Greek \emph{epsilon} and Latin \textbf{E}. (Another one for the family tree: \emph{ʾālep}$\to$A, \emph{bēt}$\to$B, \emph{hē}$\to$\textbf{E}, \emph{mem}$\to$M, \emph{ʿayin}$\to$O.)

\subsection*{Script — \ar{ة} (tāʾ marbūṭa) — the \textquotedblleft{}tied tāʾ\textquotedblright{}}
Now look carefully at this:

\begin{itemize}
\raggedright
  \item \textbf{\ar{ه}} = \emph{hāʾ}'s loop, \textbf{no dots}.
  \item \textbf{\ar{ة}} = \textbf{the same loop + tāʾ's two dots above}.
\end{itemize}

Its name says exactly what it is: \textbf{tāʾ marbūṭa}, \textquotedblleft{}the \textbf{tied-up tāʾ}\textquotedblright{} — a \textbf{\ar{ت}} (\emph{tāʾ}, Lesson 4) whose ends have been \textbf{tied into a loop}. It appears \textbf{only at the end of a word}, and it is the standard mark of a \textbf{feminine} noun or adjective.

You already know its partner: back in Lesson 4 you learned \textbf{\ar{ت}} with \textbf{two dots above}. Put those same two dots on a \emph{hāʾ} loop and you get the feminine ending. Nothing new to memorise — just two known pieces combined.

\textbf{Where you've seen it:} the farewell \textbf{\ar{مع} \ar{السلامة}} (\emph{maʿa s-salāma}) ends in \textbf{\ar{ة}} — \emph{salāma} is feminine. That's the next lesson.

\subsection*{Your turn}
\begin{itemize}
\raggedright
  \item \emph{Write it:} \textbf{\ar{ه}} — the loop; then try it isolated and final, feel it change
  \item \emph{Write it:} \textbf{\ar{ة}} — the same loop, \textbf{plus tāʾ's two dots above}
  \item \emph{Say it:} the name — \textquotedblleft{}\emph{tāʾ marbūṭa}, the \textbf{tied} tāʾ — it marks the feminine\textquotedblright{}
  \item \emph{Say it:} \textquotedblleft{}\emph{hē} $\to$ Greek epsilon $\to$ Latin \textbf{E}\textquotedblright{}
\end{itemize}

\begin{tcolorbox}[breakable,colback=teal!4,colframe=teal!35!black,title={Before you move on}]
What makes \textbf{\ar{ه}} unusual among the letters you've learned? (Its \textbf{loop changes shape} dramatically by position — four very different coats.) What is \textbf{\ar{ة}} made of, and what does its name mean? (\emph{hāʾ}'s loop \textbf{+ tāʾ's two dots}; \textquotedblleft{}\textbf{tied tāʾ}\textquotedblright{}.) What does \textbf{\ar{ة}} tell you about a word, and where can it appear? (That it is \textbf{feminine}; only at the \textbf{end}.) Next: assembling the farewell \textbf{\ar{مع} \ar{السلامة}}.
\end{tcolorbox}
\section[maʿa s-salāma]{\ar{مع} \ar{السلامة} — the goodbye, and a word you already know}

\label{lesson:AR-W12-maa-salama}



\begin{quote}
The closing move. \textbf{\ar{مع} \ar{السلامة}} (\emph{maʿa s-salāma}) is how Arabic says goodbye — literally \textquotedblleft{}\textbf{with safety},\textquotedblright{} \emph{go in peace}. And the heart of it is a word you wrote \textbf{all the way back in Lesson 3}: \textbf{\ar{سلام}}. Everything since has been leading here.
\end{quote}

\subsection*{Script — Part 1 — \ar{مع} (maʿa, \textquotedblleft{}with\textquotedblright{})}
Two letters, both known:

1. \textbf{\ar{م}} (\emph{mīm}, Lesson 3) — round head, joining left into… 2. \textbf{\ar{ع}} (\emph{ʿayn}, Lesson 10) — the open curve over its bowl, closing the word.

$\to$ \textbf{\ar{مع}} = \emph{maʿa}, \textquotedblleft{}with.\textquotedblright{}

\subsection*{Script — Part 2 — \ar{السلامة} (as-salāma, \textquotedblleft{}the safety\textquotedblright{})}
Look at what this is made of:

\begin{itemize}
\raggedright
  \item \textbf{\ar{ال}} — the \textbf{al-} (\textquotedblleft{}the\textquotedblright{}), \emph{alif} + \emph{lām}, both from Lessons 1–2. You've read it since Chapter 1 in \textbf{\ar{ال}}\ar{خير} and \textbf{\ar{ال}}\ar{سلام} \ar{عليكم}.
  \item \textbf{\ar{سلام}} — \textbf{the exact word you assembled in Lesson 3}: \emph{sīn · lām · alif · mīm}.
  \item \textbf{\ar{ة}} — the \emph{tāʾ marbūṭa} from Lesson 11, marking \emph{salāma} as \textbf{feminine}.
\end{itemize}

So \textbf{\ar{السلامة}} = \emph{al-} + \emph{salām} + \emph{-a}. You are not learning a new word; you are \textbf{adding two edges} to one you already own.

\subsection*{Script — The root s-l-m, one more time}
\emph{salām}, \emph{salāma}, \emph{islām}, \emph{muslim} — all from the Semitic root \textbf{s–l–m}, \textquotedblleft{}peace, safety, wholeness.\textquotedblright{} The greeting you learned first (\textbf{\ar{السلام} \ar{عليكم}}, \textquotedblleft{}peace be upon you\textquotedblright{}) and the farewell you learn last (\textbf{\ar{مع} \ar{السلامة}}, \textquotedblleft{}with safety\textquotedblright{}) come from the \textbf{same three consonants}. Arabic opens and closes a conversation with one idea.

\subsection*{Script — Assemble it — right to left}
1. \textbf{\ar{م}} then \textbf{\ar{ع}} $\to$ \textbf{\ar{مع}} 2. lift; \textbf{\ar{ا}} + \textbf{\ar{ل}} (the \emph{al-}), the \emph{lām} joining into… 3. \textbf{\ar{س} · \ar{ل} · \ar{ا}} — and remember the \textbf{\ar{لا}} ligature, and that \emph{alif} \textbf{breaks}… 4. \textbf{\ar{م}}, then \textbf{\ar{ة}} to close.

$\to$ \textbf{\ar{مع} \ar{السلامة}}

\subsection*{Your turn}
\begin{itemize}
\raggedright
  \item \emph{Write it:} \textbf{\ar{مع}} — mīm into ʿayn
  \item \emph{Write it:} \textbf{\ar{السلامة}} — \textquotedblleft{}the \emph{al-}, then your Lesson-3 word, then the feminine \ar{ة}\textquotedblright{}
  \item \emph{Say it:} \textquotedblleft{}\emph{maʿa s-salāma} — \textbf{with safety}\textquotedblright{}
  \item \emph{Say it:} the bookend — \textquotedblleft{}\textbf{\ar{السلام} \ar{عليكم}} to greet, \textbf{\ar{مع} \ar{السلامة}} to part — same root \textbf{s-l-m}\textquotedblright{}
\end{itemize}

\begin{tcolorbox}[breakable,colback=teal!4,colframe=teal!35!black,title={Before you move on}]
What does \textbf{\ar{مع} \ar{السلامة}} literally mean? (\textquotedblleft{}\textbf{With safety}.\textquotedblright{}) Which earlier word sits inside \textbf{\ar{السلامة}}, and what was added? (\textbf{\ar{سلام}} from Lesson 3 — plus \textbf{\ar{ال}} in front and \textbf{\ar{ة}} behind.) What does the final \textbf{\ar{ة}} signal? (\emph{salāma} is \textbf{feminine}.) What root ties the first greeting to the last farewell? (\textbf{s–l–m} — peace/safety.) You can now hand-write an Arabic conversation from hello to goodbye.
\end{tcolorbox}
\section[maʿa]{\ar{مع} (maʿa) — \textquotedblleft{}with\textquotedblright{}}

\label{lesson:AR-C04-maa-with}



\begin{quote}
Two letters, both already yours, and a word that will carry the whole chapter.
\end{quote}

\subsection*{You'll want to know — The letters}
\emph{(Skim if you read Arabic.)} \textbf{\ar{م}} (\emph{mīm}) from Lesson 3 of the writing set, and \textbf{\ar{ع}} (\emph{ʿayn}) from Lesson 10 — the throat sound whose Phoenician ancestor meant \textquotedblleft{}\textbf{eye},\textquotedblright{} and which the Greeks recycled into the vowel \textbf{O}.

Right to left: \textbf{\ar{م} · \ar{ع}} $\to$ \emph{maʿa}. The \emph{ʿayn} is not a vowel; it is a constriction deep in the throat. Take your time with it.

\subsection*{You'll want to know — The word}
\textbf{\ar{مع}} (\emph{maʿa}) = \textquotedblleft{}\textbf{with}, together with.\textquotedblright{}

It is a \textbf{preposition}, and unlike \textbf{\ar{الـ}} (\emph{al-}) or \textbf{\ar{بـ}} (\emph{bi-}) from Chapter 3, it is a \textbf{free-standing word} — it does not glue onto the next one. Compare:

\noindent
\begin{tabularx}{\linewidth}{@{}>{\raggedright\arraybackslash}X>{\raggedright\arraybackslash}X@{}}
\toprule
\textbf{} & \textbf{attaches?} \\
\midrule
\textbf{\ar{الـ}} \emph{al-} \textquotedblleft{}the\textquotedblright{} & yes \\
\textbf{\ar{بـ}} \emph{bi-} \textquotedblleft{}in/with\textquotedblright{} & yes \\
\textbf{\ar{لـ}} \emph{li-} \textquotedblleft{}to\textquotedblright{} & yes \\
\textbf{\ar{مع}} \emph{maʿa} \textquotedblleft{}with\textquotedblright{} & \textbf{no} \\
\bottomrule
\end{tabularx}

So Arabic has \emph{two} words touching on \textquotedblleft{}with\textquotedblright{}: the one-letter \textbf{\ar{بـ}} (\emph{bi-}, \textquotedblleft{}in/by means of,\textquotedblright{} as in \emph{bi-khayr}) and the full word \textbf{\ar{مع}} (\emph{maʿa}, \textquotedblleft{}accompanied by\textquotedblright{}). \emph{Bi-khayr} is \textquotedblleft{}\textbf{in} goodness\textquotedblright{}; \emph{maʿa} is standing \textbf{next to} something.

\begin{culture}
\emph{maʿa} $\leftarrow$ Semitic \emph{*ʿimma}, and its Hebrew cousin is \textbf{\he{עִם}} (\emph{ʿim}), \textquotedblleft{}with.\textquotedblright{} The same pairing you met with \textbf{salām / shalom} in Chapter 1 and \textbf{ism / shem} in Chapter 2. Three chapters, three cousins — the family resemblance is systematic, not a coincidence.
\end{culture}

\subsection*{Your turn}
\begin{itemize}
\raggedright
  \item \emph{Say it:} \textquotedblleft{}\emph{maʿa}\textquotedblright{} — the \emph{ʿayn} from the throat
  \item \emph{Say it:} the contrast — \textquotedblleft{}\emph{bi-khayr} \textbf{in} goodness · \emph{maʿa} \textbf{with}\textquotedblright{}
  \item \emph{Say it:} the cousin — \textquotedblleft{}\emph{maʿa} · Hebrew \emph{ʿim}\textquotedblright{}
  \item \emph{Write it:} \textbf{\ar{مع}} — mīm joining into ʿayn
\end{itemize}

\begin{tcolorbox}[breakable,colback=teal!4,colframe=teal!35!black,title={Before you move on}]
What does \textbf{\ar{مع}} mean? (\textquotedblleft{}\textbf{With}.\textquotedblright{}) Does it attach to the next word, like \emph{al-} and \emph{bi-}? (\textbf{No} — it stands alone.) What is the difference between \emph{bi-} and \emph{maʿa}? (\emph{Bi-} = \textquotedblleft{}\textbf{in}/by means of\textquotedblright{}; \emph{maʿa} = \textquotedblleft{}\textbf{accompanied by}\textquotedblright{}.) What is its Hebrew cousin? (\emph{\textbf{ʿim}}.) Next: the word it goes with — and you already know most of it.
\end{tcolorbox}
\section[as-salāma]{\ar{السلامة} (as-salāma) — \textquotedblleft{}the safety\textquotedblright{}}

\label{lesson:AR-C04-al-salama}



\begin{quote}
This looks like a long new word. It isn't. It is \textbf{one word you learned in Chapter 1} with a piece added at each end.
\end{quote}

\subsection*{You'll want to know — Taking it apart}
\noindent
\begin{tabularx}{\linewidth}{@{}>{\raggedright\arraybackslash}X>{\raggedright\arraybackslash}X>{\raggedright\arraybackslash}X@{}}
\toprule
\textbf{piece} & \textbf{what it is} & \textbf{met in} \\
\midrule
\textbf{\ar{الـ}} & \emph{al-}, \textquotedblleft{}the\textquotedblright{} & Ch. 1 \\
\textbf{\ar{سلام}} & \emph{salām}, \textquotedblleft{}peace\textquotedblright{} & Ch. 1 — the very first word \\
\textbf{\ar{ة}} & \emph{tāʾ marbūṭa}, feminine ending & writing Lesson 11 \\
\bottomrule
\end{tabularx}

$\to$ \textbf{\ar{السلامة}} = \emph{al-} + \emph{salām} + \emph{-a} = \textquotedblleft{}\textbf{the safety}.\textquotedblright{}

Every letter is one you have already written by hand.

\subsection*{You'll want to know — The pattern that changed the meaning}
Arabic doesn't just add endings — it pours a root into a \textbf{shape}. Here the root \textbf{s–l–m} goes into the pattern \emph{faʿāla}, which turns a thing into \textbf{the state of being that thing}:

\noindent
\begin{tabularx}{\linewidth}{@{}>{\raggedright\arraybackslash}X>{\raggedright\arraybackslash}X@{}}
\toprule
\textbf{word} & \textbf{meaning} \\
\midrule
\textbf{salām} & peace \\
\textbf{salāma} & safety, soundness, wellbeing — \emph{the state of being at peace} \\
\bottomrule
\end{tabularx}

Same three consonants, one new shape, a new but related word. This is the engine Chapter 1 promised, working in front of you.

\subsection*{You'll want to know — Two small rules, both already learned}
\textbf{1. The sun letter.} \emph{S} (\textbf{\ar{س}}) is a \textbf{sun letter}, so the \emph{l} of \emph{al-} assimilates to it: written \textbf{\ar{الس}}, but said \textbf{as-}salāma, not \emph{al-salāma}. The same thing happens in \textbf{\ar{الس}}\ar{لام} \ar{عليكم} — the greeting you have said since Chapter 1.

\textbf{2. The feminine \ar{ة}.} \emph{Tāʾ marbūṭa} marks \emph{salāma} as \textbf{feminine}. Arabic abstract nouns very often are.

\subsection*{Your turn}
\begin{itemize}
\raggedright
  \item \emph{Say it:} \textquotedblleft{}\emph{as-salāma}\textquotedblright{} — with the doubled \emph{s}, not \emph{al-s}
  \item \emph{Say it:} the build — \textquotedblleft{}\emph{salām} … \emph{salāma} … \emph{as-salāma}\textquotedblright{}
  \item \emph{Say it:} the pattern — \textquotedblleft{}peace $\to$ \textbf{the state of} peace\textquotedblright{}
  \item \emph{Write it:} \textbf{\ar{السلامة}}
\end{itemize}

\begin{tcolorbox}[breakable,colback=teal!4,colframe=teal!35!black,title={Before you move on}]
What three pieces make \textbf{\ar{السلامة}}? (\emph{\textbf{al-}} + \emph{\textbf{salām}} + \emph{\textbf{\ar{ة}}}.) What did the \emph{faʿāla} pattern do to the meaning? (Turned \textquotedblleft{}peace\textquotedblright{} into \textquotedblleft{}\textbf{the state of} being at peace\textquotedblright{} — safety.) Why is it \emph{as-salāma} and not \emph{al-salāma}? (\textbf{\ar{س}} is a \textbf{sun letter}; the \emph{l} assimilates.) What does the \textbf{\ar{ة}} tell you? (The word is \textbf{feminine}.) Next: put \emph{maʿa} in front of it.
\end{tcolorbox}
\section[maʿa s-salāma]{\ar{مع} \ar{السلامة} (maʿa s-salāma) — \textquotedblleft{}goodbye\textquotedblright{}}

\label{lesson:AR-C04-maa-salama}



\begin{quote}
Put the last two lessons together and the farewell falls out. You have already \textbf{hand-written} this phrase; now it means something.
\end{quote}

\subsection*{You'll want to know — The assembly}
\begin{quote}
\textbf{\ar{مع}} (\emph{maʿa}, \textquotedblleft{}with\textquotedblright{}) + \textbf{\ar{السلامة}} (\emph{as-salāma}, \textquotedblleft{}the safety\textquotedblright{}) $\to$ \textbf{\ar{مع} \ar{السلامة}} — \textquotedblleft{}\textbf{with safety}.\textquotedblright{}
\end{quote}

Said as one run: \emph{maʿa s-salāma} — the \emph{a} of \emph{maʿa} swallows the \emph{al-}'s vowel, and the sun-letter rule doubles the \emph{s}.

It is not \textquotedblleft{}goodbye\textquotedblright{} in the English sense of a plain sign-off. It is a \textbf{wish}: \emph{go in safety}.

\subsection*{You'll want to know — The circle closes}
Look at what Chapter 1 opened with and what Chapter 4 closes with:

\noindent
\begin{tabularx}{\linewidth}{@{}>{\raggedright\arraybackslash}X>{\raggedright\arraybackslash}X>{\raggedright\arraybackslash}X@{}}
\toprule
\textbf{} & \textbf{phrase} & \textbf{root} \\
\midrule
first greeting & \textbf{\ar{السلام} \ar{عليكم}} \emph{as-salāmu ʿalaykum} & \textbf{s–l–m} \\
last farewell & \textbf{\ar{مع} \ar{السلامة}} \emph{maʿa s-salāma} & \textbf{s–l–m} \\
\bottomrule
\end{tabularx}

\textbf{The same three consonants.} Arabic begins and ends a conversation with one idea — peace, safety, wholeness. The root engine you met in the very first lesson has been running the whole time; here it simply comes back round.

The family, one more time: \emph{salām} (peace), \emph{salāma} (safety), \emph{islām} (submission, i.e. entering peace), \emph{muslim} (one who does so).

\subsection*{You'll want to know — English does the same trick}
\textbf{Goodbye} looks like a single word. It is worn down from \emph{\textbf{God be with ye}}. So the English farewell is also built on \textquotedblleft{}\textbf{with}\textquotedblright{} — you are being commended to something protective as you go.

\noindent
\begin{tabularx}{\linewidth}{@{}>{\raggedright\arraybackslash}X>{\raggedright\arraybackslash}X>{\raggedright\arraybackslash}X@{}}
\toprule
\textbf{language} & \textbf{farewell} & \textbf{literally} \\
\midrule
Arabic & \emph{maʿa s-salāma} & \textquotedblleft{}\textbf{with} safety\textquotedblright{} \\
English & \emph{goodbye} & \textquotedblleft{}God be \textbf{with} ye\textquotedblright{} \\
Spanish & \emph{adiós} & \textquotedblleft{}to God\textquotedblright{} \\
French & \emph{adieu} & \textquotedblleft{}to God\textquotedblright{} \\
\bottomrule
\end{tabularx}

Four languages, one instinct: as someone leaves, hand them over to safekeeping.

\subsection*{Your turn}
Say these aloud:

\begin{itemize}
\raggedright
  \item \textquotedblleft{}\emph{maʿa s-salāma}\textquotedblright{}
  \item the bookend — \textquotedblleft{}\emph{as-salāmu ʿalaykum} … \emph{maʿa s-salāma}\textquotedblright{}
  \item the root aloud — \textquotedblleft{}\textbf{s · l · m}\textquotedblright{}
  \item \textquotedblleft{}goodbye $\leftarrow$ God be \textbf{with} ye\textquotedblright{}
\end{itemize}

\begin{tcolorbox}[breakable,colback=teal!4,colframe=teal!35!black,title={Before you move on}]
What does \textbf{\ar{مع} \ar{السلامة}} literally mean? (\textquotedblleft{}\textbf{With safety}.\textquotedblright{}) What root is it built on, and where have you met that root before? (\textbf{s–l–m} — the Chapter 1 greeting \emph{as-salāmu ʿalaykum}, and \emph{salām} itself.) So what does Arabic open and close a conversation with? (\textbf{The same idea} — peace/safety.) What is English \emph{goodbye} worn down from? (\textquotedblleft{}\textbf{God be with ye}.\textquotedblright{}) Next: a second farewell, and a word Spanish borrowed.
\end{tcolorbox}
\section[ilā l-liqāʾ]{\ar{إلى} \ar{اللقاء} (ilā l-liqāʾ) — \textquotedblleft{}see you\textquotedblright{}}

\label{lesson:AR-C04-ila-liqaa}



\begin{quote}
The other everyday farewell — and the one with a surprise waiting inside Spanish.
\end{quote}

\subsection*{You'll want to know — A note on the letters}
\emph{(Skim if you read Arabic.)} Two pieces here are \textbf{beyond what the writing track has covered}, so read them for now and draw them later:

\begin{itemize}
\raggedright
  \item \textbf{\ar{ق}} (\emph{qāf}) — a \emph{k} made far back against the soft palate, deeper than \textbf{\ar{ك}} (\emph{kāf}). It is \textbf{not} in the letter set you have been writing from.
  \item \textbf{\ar{ى}} in \textbf{\ar{إلى}} — this is \emph{alif maqṣūra}, \textquotedblleft{}shortened alif\textquotedblright{}: it wears \textbf{yāʾ's skeleton with no dots} and is pronounced as a long \emph{ā}. You already know from writing Lessons 4–5 that dots are what separate letters sharing a body; this is that same principle, taken one step further.
\end{itemize}

The rest — \textbf{\ar{ا} \ar{ل} \ar{ء}} — you have written.

\subsection*{You'll want to know — The two pieces}
\noindent
\begin{tabularx}{\linewidth}{@{}>{\raggedright\arraybackslash}X>{\raggedright\arraybackslash}X@{}}
\toprule
\textbf{piece} & \textbf{meaning} \\
\midrule
\textbf{\ar{إلى}} \emph{ilā} & \textquotedblleft{}to, until\textquotedblright{} \\
\textbf{\ar{اللقاء}} \emph{al-liqāʾ} & \textquotedblleft{}the meeting\textquotedblright{} \\
\bottomrule
\end{tabularx}

$\to$ \textbf{\ar{إلى} \ar{اللقاء}} = \textquotedblleft{}\textbf{until the meeting}.\textquotedblright{}

\textbf{\ar{لقاء}} (\emph{liqāʾ}) comes from the root \textbf{l–q–y}, \textquotedblleft{}\textbf{to meet}.\textquotedblright{} The final \textbf{\ar{ء}} is the \emph{hamza} from writing Lesson 6 — a glottal stop, a real consonant here, not decoration.

Note \textbf{\ar{ل}} is a sun letter too, so \emph{al-liqāʾ} doubles the \emph{l}, just as \emph{as-salāma} doubled the \emph{s}.

\begin{culture}
Spanish says goodbye like this:

\begin{quote}
\emph{hasta luego} — \textquotedblleft{}\textbf{until} later\textquotedblright{} \emph{hasta mañana} — \textquotedblleft{}\textbf{until} tomorrow\textquotedblright{}
\end{quote}

Same shape as \emph{ilā l-liqāʾ}: \textbf{\textquotedblleft{}until\textquotedblright{} + a future point}. But it goes further — Spanish \emph{\textbf{hasta}} is itself an \textbf{Arabic loanword}, from \emph{\textbf{ḥattā}} (\textquotedblleft{}until\textquotedblright{}). Of the roughly four thousand Arabic words Spanish took during the eight centuries of al-Andalus, most are nouns; \emph{hasta} is one of the rare borrowed \textbf{function words}.

So Spanish did not merely copy the pattern — it borrowed the very word that builds it.

\noindent
\begin{tabularx}{\linewidth}{@{}>{\raggedright\arraybackslash}X>{\raggedright\arraybackslash}X>{\raggedright\arraybackslash}X@{}}
\toprule
\textbf{language} & \textbf{farewell} & \textbf{shape} \\
\midrule
Arabic & \emph{ilā l-liqāʾ} & until + the meeting \\
Spanish & \emph{hasta luego} & \textbf{until} ($\leftarrow$ \emph{ḥattā}) + later \\
French & \emph{à bientôt} & to + soon \\
Italian & \emph{a presto} & to + soon \\
\bottomrule
\end{tabularx}
\end{culture}

\subsection*{Your turn}
Say these aloud:

\begin{itemize}
\raggedright
  \item \textquotedblleft{}\emph{ilā l-liqāʾ}\textquotedblright{} — double the \emph{l}
  \item the two farewells — \textquotedblleft{}\emph{maʿa s-salāma} · \emph{ilā l-liqāʾ}\textquotedblright{}
  \item the loan — \textquotedblleft{}\emph{hasta} $\leftarrow$ \emph{ḥattā}\textquotedblright{}
  \item the root — \textquotedblleft{}\textbf{l · q · y} — to meet\textquotedblright{}
\end{itemize}

\begin{tcolorbox}[breakable,colback=teal!4,colframe=teal!35!black,title={Before you move on}]
What does \textbf{\ar{إلى} \ar{اللقاء}} literally say? (\textquotedblleft{}\textbf{Until the meeting}.\textquotedblright{}) What root is \emph{liqāʾ} from? (\textbf{l–q–y}, \textquotedblleft{}to meet\textquotedblright{}.) Which two letters here has the writing track \textbf{not} reached? (\textbf{\ar{ق}} \emph{qāf}, and \textbf{\ar{ى}} \emph{alif maqṣūra}.) What is \emph{alif maqṣūra}? (\textbf{Yāʾ's body with no dots}, pronounced long \emph{ā}.) Which Spanish word for \textquotedblleft{}until\textquotedblright{} came from Arabic? (\emph{\textbf{Hasta}} $\leftarrow$ \emph{\textbf{ḥattā}}.) Next: practise the whole conversation, hello to goodbye.
\end{tcolorbox}
\section[Practice]{Practice — the whole conversation}

\label{lesson:AR-C04-practice}



\begin{quote}
Four chapters. Greet, introduce, ask after someone, part. Here it all is in one piece.
\end{quote}

\subsection*{You'll want to know — The full exchange}
\noindent
\begin{tabularx}{\linewidth}{@{}>{\raggedright\arraybackslash}X>{\raggedright\arraybackslash}X>{\raggedright\arraybackslash}X@{}}
\toprule
\textbf{} & \textbf{Arabic} & \textbf{literally} \\
\midrule
A & \textbf{\ar{السلام} \ar{عليكم}} (\emph{as-salāmu ʿalaykum}) & peace {[}be{]} upon you \\
B & \textbf{\ar{وعليكم} \ar{السلام}} (\emph{wa-ʿalaykum as-salām}) & and upon you, peace \\
A & \textbf{\ar{ما} \ar{اسمك؟}} (\emph{mā ismuka?}) & what {[}is{]} your name? \\
B & \textbf{\ar{اسمي} …\ar{،} \ar{تشرفنا}.} (\emph{ismī …, tasharrafnā}) & my name {[}is{]} …, we are honoured \\
A & \textbf{\ar{كيف} \ar{حالك؟}} (\emph{kayfa ḥāluka?}) & how {[}is{]} your state? \\
B & \textbf{\ar{الحمد} \ar{لله،} \ar{بخير}.} (\emph{al-ḥamdu lillāh, bi-khayr}) & the praise {[}is{]} to God — in goodness \\
A & \textbf{\ar{مع} \ar{السلامة}} (\emph{maʿa s-salāma}) & with safety \\
B & \textbf{\ar{إلى} \ar{اللقاء}} (\emph{ilā l-liqāʾ}) & until the meeting \\
\bottomrule
\end{tabularx}

Read the right-hand column top to bottom. \textbf{Not one line contains a verb \textquotedblleft{}to be.\textquotedblright{}} The zero copula has held for four chapters.

\begin{culture}
Everything you have learned sits on a handful of three-consonant roots:

\noindent
\begin{tabularx}{\linewidth}{@{}>{\raggedright\arraybackslash}X>{\raggedright\arraybackslash}X>{\raggedright\arraybackslash}X@{}}
\toprule
\textbf{root} & \textbf{meaning} & \textbf{words you know} \\
\midrule
\textbf{s–l–m} & peace, safety & \emph{salām}, \emph{as-salāmu ʿalaykum}, \emph{salāma}, \emph{maʿa s-salāma} \\
\textbf{s–m–w} & name & \emph{ism}, \emph{ismī}, \emph{ismuka} \\
\textbf{ḥ–w–l} & to turn & \emph{ḥāl}, \emph{kayfa ḥāluka} \\
\textbf{kh–y–r} & good & \emph{ṣabāḥ al-khayr}, \emph{bi-khayr} \\
\textbf{ḥ–m–d} & to praise & \emph{al-ḥamdu lillāh}, \emph{Muḥammad} \\
\textbf{l–q–y} & to meet & \emph{liqāʾ}, \emph{ilā l-liqāʾ} \\
\bottomrule
\end{tabularx}

Six roots. A whole conversation.
\end{culture}

\begin{grammarlens}[title={Grammar Lens — The small attaching pieces}]
\noindent
\begin{tabularx}{\linewidth}{@{}>{\raggedright\arraybackslash}X>{\raggedright\arraybackslash}X>{\raggedright\arraybackslash}X@{}}
\toprule
\textbf{piece} & \textbf{job} & \textbf{example} \\
\midrule
\textbf{\ar{الـ}} \emph{al-} & the & \emph{al-khayr} \\
\textbf{\ar{بـ}} \emph{bi-} & in / by & \emph{bi-khayr} \\
\textbf{\ar{لـ}} \emph{li-} & to & \emph{lillāh} \\
\textbf{\ar{ـي}} \emph{-ī} & my & \emph{ismī} \\
\textbf{\ar{ـك}} \emph{-ka / -ki} & your (m./f.) & \emph{ismuka, ḥāluki} \\
\bottomrule
\end{tabularx}

And \textbf{\ar{مع}} \emph{maʿa}, the one that \textbf{doesn't} attach.
\end{grammarlens}

\subsection*{Your turn}
\begin{itemize}
\raggedright
  \item \emph{Say it:} the full dialogue, both parts, addressing a \textbf{man}
  \item \emph{Say it:} it again to a \textbf{woman} — every \emph{-ka} becomes \emph{-ki}
  \item \emph{Write it:} \textbf{\ar{مع} \ar{السلامة}} — every letter is one you have drawn
  \item \emph{Say it:} the bookend — \textquotedblleft{}\textbf{\ar{السلام}} to open, \textbf{\ar{السلامة}} to close\textquotedblright{}
\end{itemize}

\begin{tcolorbox}[breakable,colback=teal!4,colframe=teal!35!black,title={Before you move on}]
Give a complete greeting-to-farewell exchange. (The dialogue above.) Which root opens \textbf{and} closes it? (\textbf{s–l–m}.) Which piece of the chapter doesn't glue onto the next word? (\textbf{\ar{مع}} \emph{maʿa}.) What single word is absent from every literal translation? (\textbf{\textquotedblleft{}Is\textquotedblright{}} — the zero copula.) Which Spanish word came from Arabic \emph{ḥattā}? (\emph{\textbf{Hasta}}.) You can now hold — and hand-write — a short Arabic conversation from beginning to end.
\end{tcolorbox}
